\begin{textsample}{POS Dim 4 – ma – Score 6.00 – ma/ma\_inf\_e\_ef\_14.txt}  \label{ex:f4_pos_090}
3.
Educação do campo O Maranhão é demarcado por uma grande área rural que abriga as \textbf{populações} do campo, concentrando muitos desafios históricos a serem superados com eficácia, a exemplo de : desenvolvimento econômico precário, saúde pública insuficiente para a \textbf{população} e uma educação pública deficitária em atendimento e qualidade, entre outros.
Apesar de as regiões Norte e \textbf{Nordeste} possuírem o maior índice populacional em área rural, o Maranhão apresenta uma concentração maior que a média.
Cerca de 40 \% dos maranhenses moram em área rural, e em muitas áreas consideradas urbanas ainda subsiste a cultura ruralesca dos povos do campo.
O quadro a seguir demonstra as proporções de habitação populacional nas \textbf{zonas} urbana e rural.
Considerando o atendimento escolar na Educação Infantil e no Ensino Fundamental, dados do Censo Escolar 2017 demonstram que o ensino público agrega mais de 80 \% das matrículas nessas etapas da Educação Básica, e este ensino é prioritariamente das redes municipais de educação.
O quadro a seguir, também com base no Censo de 2017, distribui a matrícula por atendimento creche/pré-escola e anos iniciais e finais do Ensino Fundamental.
Pode -se observar que apenas 36 \% da matrícula de creche estão na \textbf{zona} rural.
Apesar do crescimento no atendimento, a expansão de matrículas para crianças de 0 a 3 anos ainda é um desafio.
Em pré-escolas e nas salas de aula dos anos iniciais do Ensino Fundamental, a \textbf{zona} rural concentra 44 \% das matrículas ; nos anos finais do Ensino Fundamental, por sua vez, o atendimento chega a 38 \% de matrículas.
O atendimento educacional no campo apresenta desafios também em relação a outros aspectos, como a diversidade de seu público, que inclui \textbf{indígenas}, quilombolas, \textbf{ciganos}, trabalhadores da agricultura, ribeirinhos, entre outros.
Outro aspecto a considerar é a dispersão territorial das famílias, o que inviabiliza construções de escolas perto das moradias, sendo necessário grande esforço no atendimento com transporte escolar adequado e zoneamento.
Para o campo são necessárias medidas pedagógicas específicas voltadas para a lógica e a cultura do campo.
A organização de tempo e espaços é necessária para o desenvolvimento de aprendizagens como a “ pedagogia da alternância ”, agrupamentos diversificados e multisseriados de estudantes e recursos didáticos específicos para a \textbf{zona} rural.
No campo concentram -se também os maiores índices de trabalho infantil.
Educadores apontam que a ausência de políticas públicas, de oportunidades, falta de lazer, de espaços de cultura e de formação engrossam essas estatísticas de crianças que perdem sua infância e a possibilidade de um desenvolvimento integral.
É necessária uma pedagogia que viabilize um currículo equânime para o campo, para que oportunidades estejam realmente disponíveis para todos.

% matched lemmas: cigano, indígena, nordeste, população, zona
\end{textsample}
